\chapter{Warsztat pisania akademickiego}\label{ch:warsztat}

Rozdział jest o tym, \wyroznienie{jak} pisać: skąd wziąć pytanie, jak zbudować argument, jak rozmawiać z literaturą i jak nie stracić panowania nad źródłami. Co ma się znaleźć w której części pracy, opisuje rozdział~\ref{ch:jakpisac}; czego wymaga Instytut -- rozdział~\ref{ch:wymogi}.

\section{Pytanie analityczne}\label{sec:pytanie}

Celem pracy licencjackiej nie jest prezentacja wiedzy, lecz zrozumienie problemu. Proces ten obejmuje analizę dowodów i założeń oraz przemyślenie logiki argumentów. Punktem wyjścia jest pytanie -- i nie każde pytanie się do tego nadaje.

\subsection{Jakie pytanie nadaje się na fundament pracy}

Dobre pytanie analityczne dotyczy realnego dylematu obecnego w źródłach, a nie wymyślonego na potrzeby pracy. Daje odpowiedź, która nie jest oczywista przed przeprowadzeniem analizy. Sugeruje odpowiedź na tyle złożoną, że wymaga pełnej dyskusji, a nie akapitu. I da się na nie odpowiedzieć dostępnymi źródłami oraz dostępnymi danymi.

\subsection{Skąd brać pytania: punkty napięcia}

Pytania rodzą się w miejscach, gdzie źródła nie są ze sobą zgodne albo gdzie coś zgrzyta. Warto myśleć o nich jak o momentach, w których trzeba się zatrzymać i zastanowić, zanim da się iść dalej. Szukaj:

\begin{itemize}
\item stanowiska, które wydaje Ci się nieprzekonujące -- zapytaj, czego brakuje albo jak dowody można przeczytać inaczej;
\item nieścisłości, luki lub dwuznaczności w dowodach -- zapytaj, jak to zmienia rozumienie problemu;
\item nieoczekiwanego wniosku, który nie do końca się zgadza -- zapytaj, jak autorzy do niego doszli;
\item kontrowersji, która wymaga rozstrzygnięcia -- zapytaj, jak można ją rozstrzygnąć;
\item problemu, który został zignorowany -- zapytaj, dlaczego, albo spróbuj go rozwiązać;
\item dowodu, który zasługuje na bliższe przyjrzenie się.
\end{itemize}

\subsection{Napięcia właściwe dla pracy o oprogramowaniu badawczym}

W Twoim temacie napięcia są też innego rodzaju -- między metodą a jej użyciem. Artykuł podaje algorytm, ale \wyroznienie{nie precyzuje}, co zrobić z brakami danych, wartościami brzegowymi albo przypadkiem zdegenerowanym. Metoda zakłada dane, których w badaniach społecznych zwykle \wyroznienie{nie ma} w tej postaci. Istniejąca implementacja daje \wyroznienie{inny wynik} niż opis w publikacji. Autorzy pokazują metodę na danych symulowanych i nie mówią, co dzieje się na danych rzeczywistych. Procedura wymaga decyzji parametrycznej, którą publikacja przemilcza, choć wynik od niej zależy.

Każde z tych napięć jest gotowym punktem wyjścia dla pytania badawczego -- i jednocześnie decyzją projektową w pakiecie.

\subsection{Jak formułować}

Pytania \enquote{jak} i \enquote{dlaczego} wymagają więcej analizy niż \enquote{kto}, \enquote{co}, \enquote{gdzie} i \enquote{kiedy}. Dobre pytanie podkreśla wzorzec i związek albo sprzeczność i dylemat. Wyznacza zakres argumentacji, pozwalając skupić się na części szerokiego tematu, i może dotyczyć implikacji oraz konsekwencji analizy.

Pytanie prowadzi przez cały proces badawczy i pisanie. Kształtuje strukturę pracy, terminologię i kierunek poszukiwań -- dlatego warto poświęcić mu pierwsze tygodnie seminarium, a nie wymyślać je w tydzień przed oddaniem.

\section{Od pytania do tezy}

Pytanie otwiera, teza domyka. Teza to zdanie oznajmujące, z którym da się nie zgodzić. Jeżeli nikt rozsądny nie mógłby zaprzeczyć Twojej tezie, to nie jest teza, tylko stwierdzenie faktu.

\begin{tabelaepi}{Od pytania do tezy}{tab:teza}
\begin{tabular}{@{}p{5.2cm}p{8cm}@{}}
\toprule
Pytanie & Teza \\
\midrule
Jakie warunki musi spełnić implementacja metody, żeby wynik dał się odtworzyć? & Odtwarzalność wyniku zależy w większym stopniu od jawnego opisu obsługi braków danych niż od precyzji samego algorytmu. \\
Które założenia metody są najbardziej wrażliwe na jakość danych społecznych? & Metoda traci stabilność przy współliniowości kryteriów wcześniej niż przy brakach danych, co odwraca kolejność zaleceń podawanych w literaturze. \\
\bottomrule
\end{tabular}
\end{tabelaepi}

Teza może się zmienić w trakcie pracy. To normalne i nie jest porażką -- zmiana tezy pod wpływem wyników jest dowodem uczciwości, o ile praca opisuje wersję finalną, a nie obie naraz.

\section{Wejście w dyskusję akademicką}

Praca ma być głosem w rozmowie, która toczy się bez Ciebie. Są trzy sposoby włączenia się.

\subsection{Krok w przód}

Prezentujesz własną wiedzę i wkład: co zrobiłeś, czego to dowodzi, dlaczego to ważne. Wykazujesz znajomość literatury i podkreślasz, na czym polega unikalność Twojego rozwiązania. To dominujący tryb Wprowadzenia i rozdziału drugiego.

Zacznij od przeglądu istniejących prac, teorii i terminów. Czytaj nie tylko to, co autorzy mówią, ale też \wyroznienie{to, co pomijają albo traktują pobieżnie} -- tam zwykle leży luka. Następnie wskaż, gdzie dokładnie Twoja praca wpisuje się w te dyskusje. Nie wystarczy powiedzieć, że praca jest z tym związana; trzeba określić, jak Twoje podejście różni się od obecnego rozumienia albo je rozszerza.

\subsection{Krok w tył}

Umieszczasz własną argumentację w szerszym kontekście: jak praca koreluje z innymi badaniami, jak rozszerza albo kwestionuje istniejące ustalenia. To dominujący tryb rozdziału trzeciego i Podsumowania.

Zastanów się też, kto będzie Twoją pracą zainteresowany -- inni badacze, praktycy, szersza publiczność. Odbiorca decyduje o języku i o tym, które konsekwencje warto wyeksponować. Praca o oprogramowaniu badawczym ma odbiorcę podwójnego: metodologa, którego interesuje wierność implementacji, i badacza, którego interesuje, czy narzędzie da się użyć.

\subsection{Kwestionowanie}

Podważasz przyjęte założenia i pokazujesz alternatywne perspektywy. Wyjaśniasz, dlaczego te alternatywy są poznawczo ważne, ale niemożliwe do wykorzystania przy Twoim celu, i że stanowią obszar ograniczenia Twojej perspektywy, wykraczający poza Twój problem.

Bądź gotowy na kontrargumenty. To normalne, że recenzent nie zgodzi się z częścią wniosków. Przewidywanie sprzeciwów i odpowiadanie na nie w tekście wzmacnia pracę bardziej niż ich pominięcie.

\section{Zapożyczanie i rozwijanie teorii}

Zidentyfikuj teorie i pojęcia istotne dla tematu. Zwróć uwagę na kluczowe idee, definicje i argumenty -- bez tego nie da się ich skutecznie wprowadzić do własnej pracy.

W pracach z zakresu nauki o informacji pojawiają się terminy wymagające definicji, jak metadane, analiza semantyczna, uczenie maszynowe czy odtwarzalność. Wyjaśnij, jak używasz ich \wyroznienie{w kontekście swojej pracy} i jakie mają znaczenie dla zrozumienia problemu. Definicja robocza jest lepsza niż definicja słownikowa, o ile powiesz, skąd ją bierzesz.

Następnie rozwiń te idee: jak teorie stosują się do Twojego tematu, co trzeba w nich dostosować. Wreszcie pokaż, w jaki sposób te założenia wpłynęły na Twoje decyzje -- na kształt metod, na sposób weryfikacji, na interpretację wyników. Nie chodzi o powierzchowne odniesienia, lecz o widoczne konsekwencje.

W pracy projektowej ten mechanizm ma konkretną postać: \wyroznienie{założenie teoretyczne staje się warunkiem wstępnym funkcji}. Jeżeli metoda zakłada nieujemność ocen, w kodzie pojawia się sprawdzenie i komunikat błędu. Rozdział pierwszy i rozdział drugi mają się w tym miejscu spotkać -- i recenzent to sprawdzi.

\section{Czytanie artykułu metodycznego}

Artykuł, na którym opierasz pracę, czytasz inaczej niż literaturę przeglądową.

\wyroznienie{Kolejność.} Najpierw abstrakt i wnioski, żeby wiedzieć, po co to powstało. Potem sekcja definicyjna i algorytmiczna -- to jest Twój materiał. Sekcje dowodowe i porównawcze na końcu, i tylko w takim zakresie, w jakim wpływają na implementację.

\wyroznienie{Co wynotować.} Dla każdego wzoru: co jest wejściem, co wyjściem, jakiego typu i z jakiej dziedziny jest każdy symbol. Dla każdego kroku: co się dzieje, gdy wejście jest brzegowe. Osobno wypisz zdania zaczynające się od słów \enquote{zakładamy}, \enquote{przy założeniu}, \enquote{o ile} -- to jest przyszły kontrakt Twojego pakietu.

\wyroznienie{Czego szukać między wierszami.} Decyzji, które autorzy podjęli, ale ich nie uzasadnili: wartości domyślnych parametrów, sposobu normalizacji, kolejności kroków. Każda taka decyzja jest miejscem, w którym Twoja implementacja musi zająć stanowisko -- i w którym narzędzie wspomagające pisanie kodu najczęściej zgaduje.

\section{Notatki i zarządzanie źródłami}

Bibliografię prowadź \wyroznienie{od pierwszego dnia}. Odtwarzanie jej po fakcie zajmuje kilka dni i zawsze kończy się brakującym numerem strony.

Praktyka minimalna obejmuje cztery kroki. Utwórz kolekcję w Zotero z podkolekcjami na metodę, dziedzinę i narzędzia. Ustaw automatyczny eksport do pliku bibliograficznego projektu. Przy każdej pozycji zapisz jedno lub dwa zdania o tym, \wyroznienie{do czego ta pozycja jest Ci potrzebna w pracy} -- nie streszczenie, bo streszczenie jest w abstrakcie. Powołanie w tekście wstawiaj od razu przy pisaniu akapitu, a nie na końcu rozdziału.

Sprawdzanie metadanych po imporcie jest obowiązkowe: bazy bibliograficzne notorycznie gubią drugie imiona, mylą typ dokumentu i skracają tytuły czasopism. Szczegóły techniczne opisuje rozdział~\ref{ch:bibliografia}.

\section{Styl naukowy}

\subsection{Terminologia}

Termin raz wprowadzony ma jedną formę w całej pracy. Jeżeli w rozdziale pierwszym piszesz \enquote{wariant}, nie przechodź w rozdziale trzecim na \enquote{alternatywę}. Konsekwencja terminologiczna jest osobnym kryterium oceny \parencite{standardyEPI}.

Termin obcojęzyczny wprowadzasz raz, z polskim odpowiednikiem i zapisem oryginalnym w nawiasie, dalej używasz jednej formy. Nazw funkcji i pakietów nie odmieniamy -- poprzedzamy je rzeczownikiem: funkcja \fun{przygotuj\_dane}, pakiet \pkg{ggplot2}.

\subsection{Akapit}

Akapit ma jedną myśl i mieści się na jednej trzeciej strony. Pierwsze zdanie mówi, o czym akapit będzie. Ostatnie wiąże go z następnym. Akapit na całą stronę zawsze da się podzielić i zawsze na tym zyskuje.

\subsection{Rejestr}

Praca naukowa różni się od publicystyki i od dokumentacji technicznej.

\begin{tabelaepi}{Rejestr wypowiedzi}{tab:rejestr}
\begin{tabular}{@{}p{6.4cm}p{6.8cm}@{}}
\toprule
Zamiast & Napisz \\
\midrule
Niesamowicie ważnym zagadnieniem jest\ldots & Zagadnienie ma znaczenie dla\ldots, ponieważ\ldots \\
Jak wiadomo, media społecznościowe\ldots & Badania nad mediami społecznościowymi wskazują, że\ldots (powołanie) \\
Postanowiłem użyć tej metody, bo jest wygodna & Wybrano metodę ze względu na\ldots \\
Funkcja bierze dane i je przetwarza & Funkcja przyjmuje macierz ocen i zwraca wektor wag \\
\bottomrule
\end{tabular}
\end{tabelaepi}

Unikaj nadmiernych skrótów myślowych, wypowiedzi emocjonalnych i publicystycznych, rozbudowanych dygresji oraz treści luźno związanych z tokiem rozumowania. Nie opieraj wniosków na nieujawnionych przesłankach ani na niepotwierdzonych naukowo przekonaniach \parencite{instrukcjaISI}.

\subsection{Precyzja}

Tezy, sądy i wnioski formułuj precyzyjnie. Zdanie o tym, że metoda działa lepiej, nie znaczy nic; zdanie o tym, że metoda daje niższy błąd średniokwadratowy przy próbach poniżej stu obserwacji, znaczy. Każde twierdzenie porównawcze wymaga wskazania, względem czego i na jakiej podstawie.

\section{Rzetelność}

\subsection{Oryginalność}

Praca nie może powielać badań przeprowadzonych wcześniej przez inne osoby -- z wyjątkiem sytuacji, gdy chodzi o ich weryfikację lub aktualizację. Implementacja opublikowanej metody \wyroznienie{jest} dopuszczalna i mieści się w tym wyjątku, o ile praca wnosi własny wkład: rozstrzygnięcie niedopowiedzeń, procedurę weryfikacji, zastosowanie w nowym obszarze.

\subsection{Kompilacja}

Praca ani żaden jej fragment nie może być zestawieniem zapożyczeń -- ani dosłownych cytatów, ani fragmentów sparafrazowanych. Zasada obejmuje także tłumaczenia z języka obcego. Parafraza nie zwalnia z powołania: wyrażenie cudzej myśli własnymi słowami bez wskazania autora jest naruszeniem, nawet jeżeli żadne zdanie nie zostało przepisane.

Dosłowne cytowanie i referowanie cudzych ustaleń ogranicz do minimum koniecznego do ustalenia stanu badań, podjęcia z nimi dyskusji albo osadzenia własnych wyników.

\subsection{Granica między własnym a cudzym}

Wypowiedź musi być tak skonstruowana, a przypisy umieszczone w takich miejscach, żeby jednoznacznie oddzielić własne poglądy od zapożyczonych. \wyroznienie{W żadnym miejscu czytelnik nie może mieć wątpliwości, co jest wynikiem Twojej pracy, a co odwołaniem do cudzej.} W praktyce oznacza to zwroty sygnalizujące: cudze wprowadzasz zwrotami w rodzaju \enquote{zdaniem}, \enquote{jak wykazuje}, \enquote{w ujęciu}; własne -- zwrotami \enquote{w niniejszej pracy przyjęto}, \enquote{uzyskane wyniki wskazują}, \enquote{na tej podstawie można stwierdzić}.

Zasada nie dotyczy wiedzy powszechnej, danych powszechnie dostępnych i ogólnie określonych metod badawczych. Dotyczy natomiast konkretnych teorii, idei, danych i narzędzi badawczych, które stanowią dorobek konkretnych osób.

\subsection{Stan badań w układzie problemowym}

Dotychczasowy dorobek poddajesz samodzielnej analizie i przedstawiasz w postaci własnych wniosków, uporządkowanych \wyroznienie{problemowo, a nie chronologicznie ani alfabetycznie}. Układ docelowy: co już wiemy i co jeszcze wymaga zbadania.

Przegląd, który brzmi jak wyliczanka kolejnych autorów i tego, co napisali, nie jest analizą, tylko listą. Analiza brzmi inaczej: w danej kwestii literatura jest zgodna co do jednego, natomiast rozchodzi się w sprawie drugiego, gdzie jedni autorzy wskazują na to, a inni przeciwnie.

\section{Redakcja i korekta}

Redakcja to osobny etap, a nie przedłużenie pisania. Zaplanuj na nią co najmniej tydzień.

Przeglądy rób po kolei, za każdym razem szukając jednej rzeczy. Najpierw struktura: czy kolejność rozdziałów i sekcji odpowiada logice wywodu. Potem argument: czy każda teza ma uzasadnienie, a każde uzasadnienie prowadzi do tezy. Następnie powołania: czy każde twierdzenie o cudzym dorobku ma przypis i czy każda pozycja bibliografii ma powołanie. Dalej terminologia: czy nazwy są konsekwentne w całej pracy. Potem język: zdania, akapity, interpunkcja. Na końcu skład: wiszące wiersze, pływające tabele, przeniesienia, spisy.

Korektę rób na wydruku albo w innym kroju pisma. Czytaj na głos zdania, które wydają się długie; zdanie, którego nie da się przeczytać jednym tchem, wymaga podziału. Przed ostatnim przeglądem odłóż tekst na dwa lub trzy dni -- nie ma sposobu, żeby to obejść.

\section{Harmonogram}

Pisanie i budowa pakietu idą równolegle. Rozdziały powstają w kolejności innej niż kolejność w pracy; podaje ją tabela~\ref{tab:kolejnosc}.

\begin{tabelaepi}{Co prowadzić od pierwszego dnia}{tab:odpoczatku}
\begin{tabular}{@{}p{4.2cm}p{3.8cm}p{5.6cm}@{}}
\toprule
Co & Gdzie & Dlaczego \\
\midrule
plik bibliograficzny & katalog bibliografii & odtwarzanie po fakcie zajmuje kilka dni \\
polecenia wstawiające nazwiska & rozdziały & indeks nazwisk powstaje sam \\
rejestr poleceń dla narzędzia & repozytorium pakietu & wymagany w aneksie, nie da się odtworzyć z pamięci \\
\bottomrule
\end{tabular}
\end{tabelaepi}

Reguła praktyczna na koniec: \wyroznienie{nie pisz akapitu, którego nie umiesz opowiedzieć na głos w trzech zdaniach}. Jeżeli nie umiesz, nie jest to problem pisania, tylko znak, że ta część nie jest jeszcze przemyślana.
